\section{The \tool Approach}
\label{sec:approach}

%% Figure: overview of the approach (four phases + factor-family analyzers).
\begin{figure}[t]
\centering
\resizebox{\linewidth}{!}{%
\begin{tikzpicture}[
    font=\scriptsize,
    phase/.style={draw=black!55, rounded corners=2pt, fill=#1, align=center, minimum height=1.55cm,
                  text width=2.95cm, inner sep=3pt},
    art/.style={draw=black!45, fill=white, align=center, font=\scriptsize\itshape, inner sep=2pt,
                rounded corners=1pt},
    flow/.style={-{Stealth[length=4pt]}, thick, black!65},
    side/.style={-{Stealth[length=3pt]}, black!55, densely dashed},
]
  %% phases
  \node[phase=tolblue!12] (p1) {\textbf{Phase 1}\\Semantic analysis\\[1pt]
        sites $\rightarrow$ IR $\rightarrow$ factors\\alias / mutation facts};
  \node[phase=tolteal!14, right=0.75cm of p1] (p2) {\textbf{Phase 2}\\Guided generation\\[1pt]
        pairwise $\cdot$ boundary\\$A{\rightarrow}B$, $A{\rightarrow}B{\rightarrow}A$};
  \node[phase=tolorange!16, right=0.75cm of p2] (p3) {\textbf{Phase 3}\\Validation\\[1pt]
        $E_0$--$E_3$, cold / warm\\oracles O1--O7, minimize};
  \node[phase=tolmagenta!12, right=0.75cm of p3] (p4) {\textbf{Phase 4}\\Report\\[1pt]
        de-duplicate, explain\\root-cause hypothesis};
  %% artifacts between phases
  \node[art, above=0.42cm of $(p1.north)!0.5!(p2.north)$] (scs) {SCS + test obligations};
  \node[art, above=0.42cm of $(p2.north)!0.5!(p3.north)$] (tests) {contexts + sequences};
  \node[art, above=0.42cm of $(p3.north)!0.5!(p4.north)$] (fail) {confirmed failure};
  \draw[flow] (p1) -- (p2); \draw[flow] (p2) -- (p3); \draw[flow] (p3) -- (p4);
  \draw[side] (p1.north) |- (scs.west); \draw[side] (scs.east) -| (p2.north);
  \draw[side] ([xshift=3pt]p2.north) |- (tests.west); \draw[side] (tests.east) -| (p3.north);
  \draw[side] ([xshift=3pt]p3.north) |- (fail.west); \draw[side] (fail.east) -| (p4.north);
  %% inputs / outputs
  \node[art, left=0.45cm of p1, text width=1.55cm, font=\scriptsize] (in) {Python program\\+ seed input};
  \node[art, right=0.45cm of p4, text width=1.75cm, font=\scriptsize] (out) {issue draft\\+ reproducer};
  \draw[flow] (in) -- (p1); \draw[flow] (p4) -- (out);
  %% analyzers lane
  \node[draw=black!50, rounded corners=2pt, fill=black!4, below=0.55cm of $(p2.south)!0.5!(p3.south)$,
        text width=12.4cm, align=center, inner sep=4pt] (an)
        {\textbf{Factor-family analyzers} (OpInfo seeds, one factor family at a time, same oracles)\\[1pt]
         decomposition / meta $\cdot$ configuration relations $\cdot$ binding form $\cdot$ dtype promotion $\cdot$
         Python semantics\\ differentiation modes $\cdot$ optimizer step $\cdot$ export / AOTInductor $\cdot$ disk cache};
  \draw[flow] (an.north -| p3.south) -- (p3.south);
  \draw[side] (p2.south) -- (an.north -| p2.south);
\end{tikzpicture}}
\caption{Overview of \tool. Phase~1 derives the context factors a compilation site reads; Phase~2 changes one
factor at a time; Phase~3 compares eager and layered compiled executions, cold and warm; Phase~4 packages
confirmed failures. The analyzers reuse Phases~3 and~4 for factor families that are not visible in user code.}
\label{fig:overview}
\end{figure}


Figure~\ref{fig:overview} gives an overview. A compiled execution is determined by three things: the program,
its input, and the path the compiler takes. \tool derives \emph{factors} for the first two from the source and
from operator metadata, derives factors for the third from the structure of the pipeline, and executes a matrix
in which each test differs from a reference execution in exactly one factor. A fixed oracle over seven
observables decides whether the difference is a failure, and layered execution names the stage at which it
appears.

\subsection{Factor Model}
\label{sec:factors}

\paragraph{Program-side factors from the source.} \tool locates \emph{compilation sites}, the boundaries at
which Python code enters the compiler (calls and decorators of \tcompile and project-specific wrappers), and
lowers every function reachable from a site into a small dependency IR with data, control, alias, and state
edges (Figure~\ref{fig:factors}). A context attribute becomes a \emph{semantic factor} of the site when a
backward slice from a \code{Return}, \code{Raise}, \code{Mutation}, or \code{StateWrite} node reaches a read of
that attribute: the rank, shape, stride, dtype, device, contiguity, and \code{requires\_grad} of a tensor
parameter; the value of a scalar, flag, string, or container; a global or closure variable. Comparisons and
modulo tests on integer factors yield boundary values (for \code{x.shape[0] >= 32} the values 32, 31, 33);
shape--index chains are recovered through constant offsets; writes through views yield alias and mutation facts
that determine which inputs must be built as views or overlapping buffers. The factors of a
site, with the backend and compile flags, form its \emph{Semantic Context Signature} (\scs); attributes the
program never reads are pruned. The analysis is syntactic and takes about 30~ms for our 29-program corpus.

%% Figure: from a small program to its dependency graph, factors, and obligations.
\begin{figure}[t]
\centering
\begin{minipage}[c]{0.37\linewidth}
\begin{lstlisting}
def f(x, use_fast):
    n = x.shape[0]
    y = x.view(-1)
    if use_fast and n >= 32:
        y[n - 1] += 1
    return x.sum(dtype=x.dtype)
\end{lstlisting}
\end{minipage}\hfill
\begin{minipage}[c]{0.61\linewidth}
\centering
\begin{tikzpicture}[
    font=\scriptsize, x=1cm, y=0.62cm,
    attr/.style={draw=tolblue!70, fill=tolblue!10, rounded corners=1.5pt, inner sep=2pt, minimum width=1.45cm},
    ir/.style={draw=black!55, fill=white, inner sep=2pt},
    obs/.style={draw=tolorange!90!black, fill=tolorange!18, rounded corners=1.5pt, inner sep=2pt},
    d/.style={-{Stealth[length=3pt]}, black!70},
    c/.style={-{Stealth[length=3pt]}, tolred, densely dashed},
    a/.style={-{Stealth[length=3pt]}, tolteal, thick, dotted},
]
  \node[attr] (shape) at (0,3) {x.shape[0]};
  \node[attr] (flag)  at (0,2) {use\_fast};
  \node[attr] (alias) at (0,1) {alias(x, y)};
  \node[attr] (dtype) at (0,0) {x.dtype};
  \node[ir] (n)   at (1.75,3) {n};
  \node[ir] (br)  at (3.1,3) {Branch};
  \node[ir] (idx) at (3.1,2) {Index n$-$1};
  \node[ir] (mut) at (3.1,1) {Mutation y[$\cdot$]};
  \node[obs] (state) at (5.55,1) {post-state of x};
  \node[obs] (ret)   at (5.55,0) {Return};
  \draw[d] (shape) -- (n); \draw[d] (n) -- (br);
  \draw[d] (n.south) |- (idx.west);
  \draw[d] (flag.east) -- ++(0.55,0) |- ([yshift=-3pt]br.west);
  \draw[c] (br) -- (idx); \draw[d] (idx) -- (mut);
  \draw[a] (alias) -- (mut); \draw[d] (mut) -- (state);
  \draw[d] (mut.south) |- ([yshift=3pt]ret.west);
  \draw[d] (dtype) -- ([yshift=-3pt]ret.west);
\end{tikzpicture}
\end{minipage}

\vspace{3pt}
{\scriptsize
\begin{tabular}{@{}llll@{}}
\toprule
Factor & Constraint values & Oracles & Execution \\
\midrule
\code{x.shape[0]} & 32, 31, 33 (predicate); $n{-}1$ (index bound) & value, exception & cold \\
\code{use\_fast} & True, False, on each side of $n \geq 32$ & value, mutation & cold + warm \\
\code{x.dtype} & float16, float32, float64, bfloat16 & value, metadata & cold + warm \\
\code{alias(x, y)} & contiguous, non-contiguous, overlapping view & mutation, alias & cold \\
\bottomrule
\end{tabular}}
\caption{Phase~1 on a small function. Solid edges are data dependencies, dashed edges control dependencies, dotted
edges alias dependencies. Attributes that reach an observable (orange) become semantic factors; each factor yields
a test obligation. Attributes the function never reads, such as \code{x.stride}, are pruned.}
\label{fig:factors}
\end{figure}


\paragraph{Program-side factors from operator metadata.} For operator-level testing the source is a single
call, and the informative factors come from PyTorch's OpInfo database, which lists for about 700 operators the
supported dtypes, sample inputs, and reference implementations. \tool derives three factor families from it:
\emph{boundary values}, which replace sample data by infinities, NaN, the dtype extremes, values next to
singularities (\code{nextafter(1, 0)} for inverse functions, $\pi/2$, $1 \pm \epsilon$) and magnitudes near
overflow, in tensors long enough to reach vectorized code; \emph{unsupported dtypes}, which cast a sample to a
dtype the eager kernel rejects; and \emph{dtype pairs} for binary operators. Each family is a change of one
factor of one call; the operator and the shape stay fixed.

\paragraph{Path-side factors from the pipeline.} The compilation path is determined by choices the pipeline
exposes and documents as semantics-preserving. \tool treats four of them as factors.
\begin{itemize}
  \item \emph{Graph-break position.} A call to \code{torch.\_dynamo.graph\_break()} splits the captured graph
        and makes Dynamo generate a resume function that carries locals, side effects, exception state, and
        live objects across the break. Inserting the call after any statement must not change the result.
  \item \emph{Backend layer.} \code{backend="eager"} executes the captured graph with eager kernels,
        \code{"aot\_eager"} adds AOTAutograd and decompositions, \code{"inductor"} adds lowering and code
        generation. Each layer is a semantics-preserving refinement of the previous one.
  \item \emph{Execution context.} A \code{TorchDispatchMode} that only forwards calls, autocast, \code{dynamic=True},
        \code{fullgraph}, compile composition (a compiled function calling another, \code{compile(compile(f))}),
        and CUDA-graph or autotuning modes change which components run without changing what the program
        means.
  \item \emph{Artifact route.} \code{torch.export}, its serialization round trip, AOTInductor packaging and
        loading, and the lifetime of the loaded runner replace the just-in-time runtime by a deployed one.
\end{itemize}
Every path factor has a reference execution in which the factor takes its default value, and the expected
outcome is that the observables do not change. In the terminology of metamorphic testing~\cite{DBLP:journals/csur/ChenKLPTTZ18,DBLP:journals/tse/SeguraFSC16},
each path factor is a semantics-preserving transformation of the compilation, the eager run is the source
test case, and observable equality is the relation. Because the relation is the same for all factors, one
oracle serves the whole matrix. Program-side factors are transformations of the input instead of
the compilation, but they are tested against the same relation.

\subsection{One-Factor Execution Matrix}
\label{sec:matrix}

\paragraph{Seeds.} \tool does not invent programs. A seed is a valid call of a compilation site, obtained from
the function signature, from recorded calls, from reproducer scripts, from OpInfo samples, or from a corpus of
small Python programs written for the features Dynamo models itself (the \code{random} module, exceptions and
handlers, generators and their protocol, context managers, containers and their mutation, tensor subclasses;
274 programs at the time of writing). Programs in the corpus are added by mechanism: when a factor exposes a
defect in one construct, programs that exercise the same mechanism in other constructs are written before the
sweep is repeated (Section~\ref{sec:rq1}).

\paragraph{Single-factor variants.} For each obligation derived from the \scs, \tool builds contexts that
differ from the seed in exactly one program-side factor, drawing values round-robin over obligations so that a
small budget covers every factor once before any factor twice. For each path factor it builds the variant that
changes only that factor: the corpus program with a graph break after every statement, the same call under
another backend layer, inside a dispatch mode, or through the export route. A context in which eager execution
of the mutated input raises is kept when the exception is the behavior under test; programs that eager
execution rejects and the compiler accepts are a defect class of their own (Section~\ref{sec:rq3}).

%% Figure: execution matrix as a sequence diagram (cold vs warm, return trip).
\begin{figure}[t]
\centering
\begin{tikzpicture}[
    font=\scriptsize,
    lane/.style={black!45, thin},
    run/.style={draw=black!60, fill=#1, rounded corners=1.5pt, minimum width=1.25cm, minimum height=0.46cm,
                inner sep=1pt},
    cmp/.style={{Stealth[length=3pt]}-{Stealth[length=3pt]}, tolred, thick},
]
  %% lanes
  \foreach \y/\lab in {0/{eager (\EZ)}, -0.8/{cold artifact}, -1.6/{artifact compiled for $A$}} {
    \node[anchor=east] at (-0.2,\y) {\lab};
    \draw[lane] (0,\y) -- (10.4,\y);
  }
  \node[anchor=east, font=\scriptsize\itshape, black!60] at (-0.2,0.55) {time $\rightarrow$};
  %% eager runs
  \node[run=tolteal!18] (ea) at (1.0,0) {$E_A$};
  \node[run=tolteal!18] (eb) at (4.6,0) {$E_B$};
  %% cold runs
  \node[run=tolblue!14] (ca) at (1.0,-0.8) {$C_A^{\mathit{cold}}$};
  \node[run=tolblue!14] (cb) at (4.6,-0.8) {$C_B^{\mathit{cold}}$};
  %% warm sequence
  \node[run=tolblue!14] (wa) at (2.8,-1.6) {compile($A$)};
  \node[run=tolorange!22] (wb) at (6.4,-1.6) {$C_B^{\mathit{warm}}$};
  \node[run=tolorange!22] (wa2) at (8.9,-1.6) {$C_{A'}^{\mathit{warm}}$};
  \draw[-{Stealth[length=3pt]}, black!60] (wa) -- node[above, font=\tiny] {one factor changes} (wb);
  \draw[-{Stealth[length=3pt]}, black!60] (wb) -- node[above, font=\tiny] {return} (wa2);
  %% comparisons
  \draw[cmp] (eb) -- (cb);
  \draw[cmp] (cb.south east) -- (wb.north west);
  \draw[cmp] (eb.east) to[out=0, in=90] (wb.north);
  \draw[cmp] (ca.east) to[out=-10, in=160] (wa2.north west);
  %% verdict legend
  \node[anchor=west, align=left, font=\scriptsize, text width=10.2cm] at (0,-2.55)
    {\textcolor{tolred}{$\leftrightarrow$} comparisons.\;
     $E_B \neq C_B^{\mathit{cold}}$: code generation or capture.\;
     $E_B = C_B^{\mathit{cold}} \neq C_B^{\mathit{warm}}$ without recompilation: under-specialization.\;
     $C_A^{\mathit{cold}} \neq C_{A'}^{\mathit{warm}}$: stale artifact after the return trip.};
\end{tikzpicture}
\caption{Execution matrix for two contexts $A$ and $B$ that differ in one semantic factor. Recompilations are
counted exactly, so each verdict is tied to an observed cache decision.}
\label{fig:sequence}
\end{figure}


\paragraph{Context sequences.} Guards and caches are exercised by the execution matrix of
Figure~\ref{fig:sequence}: for contexts $A$ and $B$ that differ in one factor, eager and cold-compiled runs of
both, a warm run of $B$ on the artifact compiled for $A$, and the return trip. The comparison of $E_B$ with
$C_B^{\mathit{cold}}$ and with $C_B^{\mathit{warm}}$ separates code generation from under-specialization, and
the return trip exposes incorrect invalidation. Recompilations are counted through Dynamo's frame counters with
disk caches disabled.

\subsection{Oracle}
\label{sec:oracle}

\begin{table}[t]
\caption{The seven observables compared by \tool and the mechanisms each one exposes. Cold and warm executions
of Figure~\ref{fig:sequence} are compared on every observable.}
\label{tab:oracles}
\footnotesize
\begin{tabular}{@{}lp{5.1cm}p{6.0cm}@{}}
\toprule
Observable & Compared observation & Typical mechanism exposed \\
\midrule
O1 Value & scalars, tensors, nested containers; float64 or 30-digit reference with a dtype-aware noise floor; NaN/Inf pattern; sign of zero & wrong generated loop, unstable fused reduction, fabricated value on an empty reduction, vector polynomial with absolute-only accuracy \\
O2 Metadata & shape, dtype, stride, device, layout & wrong dtype promotion in a lowering, lost \code{channels\_last}, ignored \code{device=} \\
O3 Identity & output \emph{is} an input; storage sharing between inputs and outputs & no-op elimination that returns the input object, view replaced by a copy, generator replaced by an iterator \\
O4 Exception & raised or not; innermost exception type; whether the program's handler ran & missing validation, internal assertion, compile-only failure, exception that bypasses \code{except} \\
O5 State & post-state of inputs, module buffers, reachable Python objects; count of external effects & dropped or reordered in-place update, dropped container mutation, side effect replayed twice \\
O6 Gradient & input and parameter gradients against a higher-precision reference & wrong backward after tracing, lost in-place check, zero gradient from underflow \\
O7 Process & exit code and signal of the child process & native crash, hang, runtime destroyed while its threads run \\
\bottomrule
\end{tabular}
\end{table}


Table~\ref{tab:oracles} lists the seven observables. Values and gradients are compared against a reference,
not against a fixed tolerance: every test is also executed in eager mode with floating-point inputs promoted to
float64, and for elementary functions against a 30-digit multi-precision library, and a compiled result is
flagged only if its distance to the reference exceeds the distance of the eager result by a dtype-dependent
factor, or if its NaN/Inf pattern or the sign of a zero differs. This rule accepts compiled results that are
more accurate than eager execution and rejects one-ULP differences of re-associated reductions. Metadata covers
dtype, shape, stride, device, and layout. Identity records whether an output \emph{is} an input or shares its
storage, because a compiler that returns the input object where eager execution returns a copy changes the
meaning of every later in-place update. The exception observable compares the innermost exception of the
\code{\_\_cause\_\_} chain and, when the program contains a handler, whether the handler ran. State covers the
post-state of inputs, module buffers, and Python objects reachable from the program (lists, dictionaries,
generators), and the count of external effects such as printed lines. The process observable records exit codes
and signals, since several findings terminate the interpreter.

\subsection{Localization, Confirmation, and Reports}
\label{sec:phase3}

A failing test is re-executed on the backend layers \EO--\ER (and on the Triton path when a GPU is present).
The first layer that disagrees with its predecessor names the suspected stage: graph capture, functionalization
and decomposition, or code generation; a disagreement that appears only in a warm run names specialization or
cache reuse; a disagreement that also appears without \tcompile, for example under a forwarding dispatch mode
or in the eager kernel that the compiled path calls, is attributed \emph{outside} the pipeline and reported as
such. A candidate must survive a deterministic rerun with a fixed seed, a check that eager execution is
deterministic on the input, and a rerun in a fresh interpreter; every work item of a sweep runs in a child
process so that native crashes become findings instead of ending the sweep. Inputs are then minimized over
rank, shape, values, container sizes, flags, and arguments, and the program is minimized by deleting statements
outside the static slice of the triggering factor. A confirmed failure is emitted with the minimal reproducer, the observations on every layer, the factor that
triggered it, and, where the generated code allows it, the responsible source location and a candidate patch.
Likely duplicates are annotated, never filed; submission remains a human decision.
